% Proceso de calibración del optimizador, de los datos a la evaluación en las listas de prueba.
% Misma paleta que c4-contenedores.tex. Las cifras salen de los macros \cal* de sections/cifras_calibracion.tex.

\definecolor{tkPerF}{HTML}{FFF0B8}\definecolor{tkPerB}{HTML}{9A7B00}
\definecolor{tkAppF}{HTML}{CFE2F7}\definecolor{tkAppB}{HTML}{1F4E86}
\definecolor{tkProcF}{HTML}{D3EDD0}\definecolor{tkProcB}{HTML}{2E6B2B}
\definecolor{tkExtF}{HTML}{FAD7B5}\definecolor{tkExtB}{HTML}{A8500F}
\definecolor{tkInfF}{HTML}{DEDEDE}\definecolor{tkInfB}{HTML}{4D4D4D}
\definecolor{tkBdF}{HTML}{CDEBEB}\definecolor{tkBdB}{HTML}{1C6E6E}
\definecolor{tkMqF}{HTML}{E2D6F4}\definecolor{tkMqB}{HTML}{5B2E99}
\begin{tikzpicture}[x=1cm, y=1cm, font=\scriptsize,
  per/.style={draw, thick, rounded corners=7pt, align=center, draw=tkPerB, fill=tkPerF, inner sep=3pt, minimum height=7mm},
  cont/.style={draw, rounded corners=2pt, align=center, draw=tkAppB, fill=tkAppF, inner sep=2pt, minimum height=8mm},
  proc/.style={draw, rounded corners=2pt, align=center, draw=tkProcB, fill=tkProcF, inner sep=2pt, minimum height=7mm},
  ext/.style={draw, rounded corners=2pt, align=center, draw=tkExtB, fill=tkExtF, inner sep=2pt, minimum height=7mm},
  infra/.style={draw, rounded corners=2pt, align=center, draw=tkInfB, fill=tkInfF, inner sep=2pt, minimum height=7mm},
  bd/.style={draw, cylinder, shape border rotate=90, aspect=0.22, align=center, fill=tkBdF,
    minimum width=1.25cm, draw=tkBdB, minimum height=8mm, inner sep=1pt},
  grupo/.style={draw=black!55, rounded corners=3pt, fill=black!4},
  http/.style={-{Stealth[length=1.6mm]}, semithick},
  evt/.style={-{Stealth[length=1.6mm]}, densely dashed, semithick},
  cola/.style={-{Stealth[length=1.6mm]}, semithick, draw=tkMqB},
  hueco/.style={draw=white, line width=2.8pt},
  paso/.style={circle, draw, fill=white, inner sep=0.5pt, font=\scriptsize\bfseries, minimum size=4.4mm},
  rot/.style={font=\scriptsize, fill=white, inner sep=1pt},
  etapa/.style={draw, thick, rounded corners=3pt, align=left, inner sep=4pt, text width=4.7cm, minimum height=3.1cm, anchor=north west},
  tit/.style={font=\footnotesize\bfseries, anchor=north west, inner sep=0pt},
  ref/.style={font=\scriptsize\itshape, anchor=south west, inner sep=0pt}]
\def\caja#1#2#3#4#5#6{% color, posición, título, cuerpo, dónde se ve, nombre
  \node[etapa, draw=tk#1B, fill=tk#1F] (#6) at #2 {};
  \node[tit] at ([shift={(4pt,-4pt)}]#6.north west) {#3};
  \node[anchor=north west, text width=4.4cm, align=left, inner sep=0pt] at ([shift={(4pt,-0.75cm)}]#6.north west) {#4};
  \node[ref] at ([shift={(4pt,4pt)}]#6.south west) {#5};}
\caja{Bd}{(0,0)}{1. Datos y pliegues}{\calListas{} listas de \calColegios{} colegios.\\
  Prueba: \calColegiosPrueba{} colegios, \calHVListas{} listas, apartadas antes de mirar resultados.\\
  Ajuste: \calPliegues{} pliegues por colegio, \calCVListas{} listas.}{Tabla~\ref{tab:calopt-rejilla}}{p1}
\caja{App}{(5.7,0)}{2. Umbral: exacto o NSGA-II}{\calSubproblemas{} subproblemas medidos.\\
  Ajuste del exacto: $t\approx\calConstante\,S^{\calExponente}$ s.\\
  Cruce de los dos métodos y umbral $\Lambda=\calLimite$.}{Fig.~\ref{fig:calibracion-graficas}a, Tab.~\ref{tab:calopt-umbral}}{p2}
\caja{App}{(11.4,0)}{3. Rejilla $a$, $b$}{$a,b\in\{\calFactores\}$.\\
  Se elige la configuración más barata a menos de \calUmbralHV{} del mejor hipervolumen.\\
  Elegida: $a=b=\calA$.}{Fig.~\ref{fig:calibracion-graficas}b, Tab.~\ref{tab:calopt-rejilla}}{p3}
\caja{Mq}{(0,-4.3)}{4. Pesos MAUT desde la P21}{\calEncuestados{} acudientes ordenan cuatro criterios.\\
  Precio \calWPrecio, marca \calWMarca, calidad \calWCalidad, personaje \calWPersonaje.\\
  IC 95\,\% por \emph{bootstrap}.}{Fig.~\ref{fig:calibracion-graficas}c}{p4}
\caja{Mq}{(5.7,-4.3)}{5. Sensibilidad Dirichlet}{\calSensListas{} listas y \calSensPerturbaciones{} sorteos de pesos dentro de los intervalos.\\
  El kit recomendado se mantiene en al menos \calSensICMin{} de los sorteos.}{Fig.~\ref{fig:calibracion-graficas}d, Tab.~\ref{tab:calopt-sensibilidad}}{p5}
\caja{Proc}{(11.4,-4.3)}{6. Umbral de similitud $\tau$}{$\tau=0{,}75$ se evalúa y se conserva.\\
  \calTauPares{} pares de \calTauColegios{} colegios; AUC \calTauAUC, precisión \calTauPrecision.}{Fig.~\ref{fig:tau-clasificacion}, Sec.~\ref{sec:recuperacion}}{p6}
\caja{Ext}{(0,-8.6)}{7. Evaluación en prueba}{Hipervolumen relativo al frente exacto: \calHVFrente{} en las \calHVListas{} listas; \calPruebaElegidaHV{} en las \calHVListasNsga{} con NSGA-II.\\
  Un solo kit cubre menos: MAUT \calHVMAUTPesos, Ahorrador \calHVAhorrador, Premium \calHVPremium.}{Fig.~\ref{fig:calibracion-graficas}e, Tab.~\ref{tab:calopt-hv}}{p7}
% Flechas del proceso
\draw[http] (p1.east |- 0,-1.5) -- (p2.west |- 0,-1.5);
\draw[http] (p2.east |- 0,-1.5) -- (p3.west |- 0,-1.5);
\draw[http] (p3.south) -- ++(0,-0.55) -| node[rot, pos=0.25] {configuración elegida} (p4.north);
\draw[http] (p4.east |- 0,-5.8) -- (p5.west |- 0,-5.8);
\draw[http] (p5.east |- 0,-5.8) -- (p6.west |- 0,-5.8);
\draw[http] (p6.south) -- ++(0,-0.55) -| node[rot, pos=0.25] {pesos y $\tau$ fijados} (p7.north);
% Leyenda
\def\tkley#1#2#3#4{\node[#1, minimum width=7mm, minimum height=4.2mm, inner sep=0pt, text width=7mm] at #2 {};
  \node[anchor=west, font=\scriptsize, inner sep=1pt] at ([xshift=4.5mm]#2) {#3};}
\node[draw=black!40, rounded corners=2pt, anchor=north west, minimum width=10.2cm, minimum height=1.9cm,
  inner sep=0pt] (cajaley) at (6.2,-8.6) {};
\node[anchor=north west, font=\scriptsize\bfseries] at ([shift={(2mm,-1mm)}]cajaley.north west) {Leyenda};
\begin{scope}[shift={(cajaley.north west)}]
\node[draw=tkBdB, fill=tkBdF, rounded corners=2pt, minimum width=7mm, minimum height=4.2mm] at (0.9,-0.75) {};
\node[anchor=west, font=\scriptsize] at (1.4,-0.75) {Datos};
\node[draw=tkAppB, fill=tkAppF, rounded corners=2pt, minimum width=7mm, minimum height=4.2mm] at (3.2,-0.75) {};
\node[anchor=west, font=\scriptsize] at (3.7,-0.75) {Optimizador};
\node[draw=tkMqB, fill=tkMqF, rounded corners=2pt, minimum width=7mm, minimum height=4.2mm] at (6.2,-0.75) {};
\node[anchor=west, font=\scriptsize] at (6.7,-0.75) {Pesos MAUT};
\node[draw=tkProcB, fill=tkProcF, rounded corners=2pt, minimum width=7mm, minimum height=4.2mm] at (0.9,-1.4) {};
\node[anchor=west, font=\scriptsize] at (1.4,-1.4) {Umbral $\tau$};
\node[draw=tkExtB, fill=tkExtF, rounded corners=2pt, minimum width=7mm, minimum height=4.2mm] at (3.2,-1.4) {};
\node[anchor=west, font=\scriptsize] at (3.7,-1.4) {Evaluación};
\draw[http] (6.0,-1.4) -- +(0.7,0); \node[anchor=west, font=\scriptsize] at (6.8,-1.4) {Paso siguiente};
\end{scope}
\end{tikzpicture}
